%%%%%%%%%%%%%%%%%%%%%%%%%%%%%%%%%%%%%%%%%%%%%%%%%%%%%%%%%%%%%%%%%%%%%%%%%%%%%%%%
% Architectural Reconstruction and Controlled Mutation:
% A Unified Framework for Advanced Large Language Model Safety Defenses
%
% Format: USENIX (usenix-2020-09.sty)
% Perspective: Cross-disciplinary Professor \& Anthropic/Google Chief Security Architect
%%%%%%%%%%%%%%%%%%%%%%%%%%%%%%%%%%%%%%%%%%%%%%%%%%%%%%%%%%%%%%%%%%%%%%%%%%%%%%%%

\documentclass[letterpaper,twocolumn,10pt]{article}
\usepackage{usenix-2020-09}

\usepackage{amsmath,amssymb,amsthm}
\usepackage{algorithm}
\usepackage{algpseudocode}
\usepackage{booktabs}
\usepackage{multirow}
\usepackage{xspace}

%-------------------------------------------------------------------------------
\title{\Large \bf Architectural Reconstruction and Controlled\\
  Mutation: A Unified Framework for Advanced Large\\
  Language Model Safety Defenses}

\author{
{\rm Author One}\\
Anthropic \and
{\rm Author Two}\\
Google
\thanks{Written from the perspective of cross-disciplinary professorship and Anthropic/Google chief security architecture. Substitute author names as needed.}
}

\date{}

%-------------------------------------------------------------------------------
\begin{document}
%-------------------------------------------------------------------------------
\maketitle

%-------------------------------------------------------------------------------
\begin{abstract}
The transition of Large Language Models (LLMs) from research tools to core components in critical systems has strained traditional safety paradigms. We present a unified framework that couples \emph{architectural reconstruction}---runtime adaptation via self-improving designs, context-construction defenses, and a hybrid judge ensemble---with \emph{controlled mutation} anchored to meta-cognitive monitoring and exemplar-guided silver pipelines. We formalize the interplay and \emph{decouple} safety from fragile LLM self-reflection: adaptation is driven by context confusion (decoys, ID alignment), golden exemplars, and a calibrated judge network. We synthesize SISF, context-construction defenses (ID alignment, decoy injection, self-reminder), hybrid judge ensemble with distributional evaluation and mandatory human calibration, epistemic grounding, meta-cognitive monitor and exemplar-guided mutation, and slot-based and silver-dataset pipelines into a single, hostile-examiner-hardened architecture. We discuss threat models (including task reframing, psychological manipulation, and AI judge bias), metrics, and governance, and argue that this deep refactoring is necessary for scalable, auditable LLM safety at machine speed.
\end{abstract}

%-------------------------------------------------------------------------------
\section{Introduction}
%-------------------------------------------------------------------------------

Large Language Models have ceased to be peripheral research artifacts and now underpin mission-critical infrastructure: code generation, agents with system access, negotiations, and compliance-sensitive workflows. This shift exposes a fundamental mismatch: conventional software engineering and safety assurance assume deterministic, exhaustively specifiable components, whereas LLM-driven systems exhibit emergent, stochastic behavior that is not fully specified at design time~\cite{sisf-arxiv,agentos-arxiv}. Static defenses and periodic human red-teaming operate at human speed and cannot keep pace with automated adversarial generation; when a novel jailbreak or prompt injection succeeds, the architecture often remains unchanged and remains vulnerable to the same class of attack.

Two conceptual pillars address this gap. \textbf{Architectural reconstruction} denotes the autonomous, runtime adaptation of a system's safety layer---moving safety from a one-time pre-deployment audit to a continuous, AI-driven process that synthesizes and deploys new policies in response to detected failures~\cite{sisf-arxiv}. \textbf{Controlled mutation} denotes the rigorous, semantics-preserving perturbation of prompts, demonstrations, and benchmarks to diagnose robustness, transfer expert behavior into agents, and expand evaluation coverage without test-set contamination~\cite{reflective-mutation,adm-es-arxiv}. Neither alone suffices: reconstruction without a principled mutation substrate risks overfitting to a narrow threat model; mutation without architectural adaptation leaves defenses static after each evaluation cycle.

\textbf{Safety decoupling.} The framework explicitly recognizes that LLM cognition and reflection are fragile under adversarial pressure (task reframing, psychological manipulation, judge bias). We therefore \emph{decouple} safety from the model's unfettered self-improvement. Adaptation is not the LLM ``stepping on its own foot'' blindly; it is channeled through \emph{context confusion} (decoys and ID alignment) to deny the attacker a clean operating space, through \emph{golden exemplars} to anchor defender values, and through a \emph{diverse judge ensemble} with mandatory human calibration to ensure the evolution direction remains correct.

This paper unifies both under a single framework. We assume the dual role of cross-disciplinary academic and industry chief security architect (Anthropic \& Google) to bridge theory and deployment. Our contributions are:

\begin{itemize}
\item A \emph{unified conceptual model} linking architectural reconstruction and controlled mutation, with clear boundaries (assets, oracles, feedback loops) and threat models (\S\ref{sec:framework}).
\item A \emph{reconstruction axis}: SISF, context-construction defenses (ID alignment, decoy injection, self-reminder), hybrid judge ensemble with distributional evaluation and mandatory human calibration, and epistemic grounding (\S\ref{sec:reconstruction}).
\item A \emph{mutation axis}: meta-cognitive monitor and adversarial reset, exemplar-guided mutation (golden dataset and hierarchical experience memory), slot-based mutation, silver pipelines (ADM-ES), and metadata-layer controls (\S\ref{sec:mutation}).
\item \emph{Integration, metrics, and governance}: how reconstruction and mutation feed each other, evaluation protocol, and implications for attestation and release gating (\S\ref{sec:integration}).
\end{itemize}

%-------------------------------------------------------------------------------
\section{Unified Framework and Threat Model}
\label{sec:framework}
%-------------------------------------------------------------------------------

\subsection{Assets and Oracles}

The \emph{asset} is the effective ``safety surface'' of the system: policies, guardrails, trigger sets, and the data used to evaluate and tune them. The \emph{oracle} is the process that labels behavior (safe/unsafe, compliant/non-compliant). In multi-round settings, the evaluator (or adversary) interacts with the system and observes outcomes---binary decisions, scores, or failure reports---forming a statistical query interface. Uncontrolled exposure of the same items across rounds enables reconstruction attacks that recover trigger patterns or benchmark content~\cite{llm-reconstruction-emergent,ControlledMutation_SafetyEval}.

\subsection{Reconstruction and Mutation as Dual Levers}

\textbf{Architectural reconstruction} changes \emph{who} enforces safety and \emph{when}: from fixed, human-authored rules to adaptive modules (adjudicator, policy synthesis, feedback loop) that update at runtime. It does not by itself control \emph{what} data is exposed to the oracle. \textbf{Controlled mutation} changes \emph{what} is exposed: instead of serving canonical items repeatedly, the system serves semantic-equivalent variants from a mutator, with round scheduling and dispersion to prevent gradient-based or optimization-based inversion~\cite{ControlledMutation_SafetyEval}. Together, reconstruction raises the ceiling of adaptability; mutation protects the asset and keeps the oracle informative.

\subsection{Threat Model and Hostile-Examiner Attack Surfaces}

We consider (i) \emph{external} threats: multi-round extraction of safety data, jailbreaks, and prompt-injection exfiltration; (ii) \emph{internal} threats: synthetic pollution, bias in auto-labeled or mutated data, and epistemic manipulation~\cite{epistemic-grounding}. For \emph{hostile-examiner} red teams, three attack surfaces demand deep architectural (not merely parametric) response:

\textbf{Task reframing:} The adversary disguises extraction as a legitimate context-processing task so intent-based filters see only ``normal'' behavior. Defenses must contaminate the attacker's operating space (context-construction defenses, \S\ref{sec:reconstruction}).

\textbf{Psychological manipulation (HPM):} Gaslighting and cognitive-pressure attacks induce self-doubt and value drift in the model; reflection and self-diagnosis can then produce \emph{malicious} variants. Defenses must anchor mutation to external cognitive anchors and independent monitoring (\S\ref{sec:mutation}).

\textbf{AI judge bias:} A single dynamic judge has systematic blind spots; feedback can steer the loop in the wrong direction. Defenses must use a diverse judge ensemble and mandatory human calibration (\S\ref{sec:judge}). Defenses must preserve \emph{evaluation fidelity} and \emph{auditability}.

%-------------------------------------------------------------------------------
\section{Architectural Reconstruction}
\label{sec:reconstruction}
%-------------------------------------------------------------------------------

\subsection{Self-Improving Safety Frameworks (SISF)}

SISF reframes safety as a continuous process~\cite{sisf-arxiv}. An unaligned base LLM provides capability; breach detection is delegated not to a single adjudicator but to a \emph{hybrid judge ensemble} (\S\ref{sec:judge}); a policy synthesis module consumes consensus signals and generates new heuristic and semantic policies; a dynamic feedback loop deploys them immediately. Thus the \emph{architecture} reconstructs itself in response to failure. On AdvBench-style evaluations, such frameworks can reduce Attack Success Rate from 100\% to roughly 45\% while keeping False Positive Rate at 0\% on benign inputs when the adjudication pipeline is properly calibrated~\cite{sisf-arxiv}. The key insight: when a novel attack succeeds, the system does not remain static---it synthesizes and applies a new defense within the same operational loop.

\subsection{Context-Construction Defenses (Beyond Intent Vetting)}

Intent-based ``semantic firewalls''~\cite{agentos-arxiv} fail when adversaries use \emph{task reframing}: extraction is disguised as legitimate context-processing tasks, so malicious intent is no longer detectable. We upgrade to \textbf{context-construction defenses} that do not merely filter input but \emph{actively contaminate the attacker's operating space}.

\textbf{ID Alignment:} Before retrieved context or intermediate results reach the unaligned base LLM, an \textbf{ID Alignment} module strips or normalizes all instance identifiers. Attackers can no longer refer to specific IDs to precisely scope and extract sensitive spans; benign Q\&A is unaffected. This reduces reconstruction precision without changing nominal task semantics.

\textbf{Dynamic decoy injection and risk scoring:} The system injects \emph{decoy fields} (synthetic entities or associations) with \textbf{cryptographic watermarking} (e.g., hash-derived structure) so that accidental hallucination reproducing them is extremely unlikely. Backend instructions mark decoys as internally forbidden. Benign users ignore them; task-reframing attackers often override instructions and surface decoy data. To avoid unacceptable FPR from ``one-strike'' termination (e.g., long-context agents where attention drift or harmless summarization might briefly touch a decoy), we use a \textbf{decoy-hit density} and \textbf{risk score}: a single low-fidelity touch triggers a warning and intent re-verification (e.g., re-prompt for clarification); only when the model attempts \emph{high-fidelity extraction of the full decoy structure} does the system terminate the session and log the event for audit. This keeps availability for legitimate long-context workloads while still catching deliberate extraction.

\textbf{Self-reminder at final turn:} In the last turn of multi-turn interaction, safety instructions are forcibly appended to the user query as a system prompt, reinforcing boundary awareness at generation time and reducing end-of-session drift.

\subsection{Hybrid Judge Ensemble and Calibration}
\label{sec:judge}

A single AI adjudicator (e.g., GPT-4o) exhibits systematic bias and blind spots for certain attack types and victim models. We replace it with a \textbf{hybrid judge ensemble}: multiple judges with different architectures and training distributions (e.g., Llama-Guard~\cite{llama-guard}, StrongREJECT~\cite{strongreject}, or task-specific classifiers). Only when judges reach \emph{consensus} is the feedback signal passed to the policy-synthesis module. This reduces single-model bias and prevents one compromised or skewed judge from driving the loop.

Judgment is not based on a single greedy decode. The pipeline uses \textbf{distributional evaluation}: each variant is sampled multiple times, and the system estimates the \emph{empirical refusal-failure probability distribution}, countering the inherent randomness of safety refusals. For ambiguous boundary cases, \textbf{epistemic grounding} (\S\ref{sec:epistemic}) acts as the ensemble's \emph{core weighting principle}: judge outputs that exhibit higher \emph{logical consistency} (via recursive consistency checks) receive higher confidence weight than those driven mainly by fluency or user conformity, so consensus is biased toward truth-seeking rather than hollow compliance. A \textbf{mandatory human-calibration} branch is wired in: a fixed proportion (e.g., 5\%) of high-confidence policy or variant decisions is randomly selected for human audit. This quantifies judge-ensemble FPR and corrects systematic bias, closing the loop on AI-judge reliability.

\subsection{Epistemic Grounding}
\label{sec:epistemic}

RLHF and similar alignment can prioritize fluency and user conformity over truth-seeking, yielding ``epistemic hollowing''~\cite{epistemic-grounding}. In strategic machine-to-machine settings (e.g., negotiations), adversarial agents can exploit this, causing measurable economic waste. Epistemic grounding reconstructs validation loops: model tiering (stronger model verifies weaker), recursive consistency checks, and risk-tiered handling (low/med/high) for manipulation-prone interactions. It is \emph{integrated into the hybrid judge network} (\S\ref{sec:judge}): when the ensemble aggregates votes on borderline cases, it uses consistency-check outcomes to weight each judge's contribution---judges that pass recursive validation receive higher weight than those that merely favor fluent or conforming outputs. This completes the reconstruction axis: not only \emph{what} the system does (SISF, context defenses, judge ensemble) but \emph{how} it validates its own beliefs and how that validation is wired into consensus.

%-------------------------------------------------------------------------------
\section{Controlled Mutation}
\label{sec:mutation}
%-------------------------------------------------------------------------------

\subsection{Meta-Cognitive Monitor and Adversarial Reset}
\label{sec:metacog}

Under psychological manipulation (e.g., gaslighting), the model can enter cognitive dissonance and self-doubt; ``free'' reflection may then produce unsafe or value-drifted variants. We therefore do \emph{not} grant the LLM full autonomy in self-diagnosis. The \textbf{meta-cognitive monitor} is \emph{not} another full-scale dialogue LLM (which would remain vulnerable to indirect HPM). It is implemented as either (i) a \textbf{lightweight dedicated classifier} (e.g., a small language model, SLM) trained only on stability/breach labels, or (ii) a \textbf{hidden-state activation monitor} operating at white-box level. In the latter, \textbf{representation engineering}~\cite{rep-eng} is used: the system reads activation distributions at specified layers (e.g., middle layers) of the main model to detect feature vectors associated with ``deception,'' ``jailbreak compliance,'' or ``self-cognition collapse,'' rather than relying on surface text output. Thus the monitor cannot be fooled by natural-language-level manipulation; it operates on internal representations. When it detects strategy-state or persona-parameter shifts indicative of self-doubt or cognitive instability, it triggers an \textbf{adversarial reset}: the session is sanitized, mutation state is re-anchored to the golden exemplar set, and the loop continues from a known-good configuration.

\subsection{Exemplar-Guided Mutation (No Free Reflection)}

We replace unconstrained ``reflective'' mutation~\cite{reflective-mutation} with \textbf{exemplar-guided mutation}. The mutation module is driven by a human-audited \textbf{golden dataset} and a \textbf{hierarchical experience memory} of past success/failure items. On diagnosed failure, the system retrieves semantically similar \emph{experience items} from global memory and constrains the model to produce variants conforming to \textbf{silver-standard mutations}. To keep \emph{end-to-end latency} and \emph{token cost} acceptable for real-time adversarial settings, retrieval is implemented via a lightweight \textbf{vector database} (e.g., Faiss-based) with \textbf{Top-$K$} retrieval (e.g., $K=2$) of the most relevant silver-standard mutation pairs; optional attention caching can reuse prior context when the same exemplar neighborhood is revisited. Only these $K$ exemplars are injected into the context window, avoiding unbounded context growth. Correction remains anchored to defender-defined exemplars; evolutionary dynamics still apply~\cite{metaheuristic-mdpi}, with the mutation operator bound to the exemplar set. Variants remain \emph{semantics-preserving} on the oracle~\cite{ControlledMutation_SafetyEval,jade-linguistics}.

\subsection{Silver Pipelines and Expert Behavior Transfer (ADM-ES)}

Transferring expert behavior into agents requires diagnosing latent failures (e.g., phrasing bias, extraction drift) that end-task metrics miss~\cite{adm-es-arxiv}. ADM-ES builds ``silver'' datasets via controlled mutation of golden expert data: retrieve exemplars, define mutation objectives (tone, constraint), run an LLM mutator, and accept variants under a quality check (e.g., BERTScore + consistency). A fitness score combines semantic similarity and behavioral consistency. The resulting diagnostics feed into a recommendation map (e.g., UMAP) to identify failure clusters and prescribe targeted fixes. Thus mutation is both \emph{evaluative} and \emph{remedial}.

\subsection{Slot-Based Mutation and CSE}

For complex behaviors (e.g., code optimization), \emph{slot-based} decomposition targets specific functional components rather than mutating globally~\cite{cse-evocontrol,cse-arxiv}. Controlled Self-Evolution (CSE) uses diversified planning, slot-based mutation, and compositional crossover; it outperforms base LLMs on efficiency benchmarks (e.g., EffiBench-X) and avoids plateaus typical of naive genetic search. In safety terms, slot-based mutation allows surgical stress-testing of sub-capabilities (e.g., tool routing, refusal consistency) without destroying overall semantics.

\subsection{Cognitive Taxonomy and Dynamic Benchmarking}

Static benchmarks are prone to contamination and memorization. Cognitive taxonomies (e.g., Bloom: Remember, Understand, Apply, Analyze, Evaluate, Create) structure evaluation by capability layer~\cite{bloom-bugreport}. Controlled mutation extends benchmarks through linguistic and contextual variation, forcing models to demonstrate reasoning rather than recall. In mathematical domains, ``forging'' via parameter and constraint mutation expands benchmarks (e.g., AIMO) while preserving well-posedness and auditability~\cite{aimo-benchmark}.

\subsection{Metadata and Architectural Consistency}

Safety also depends on \emph{metadata}: tool descriptions (MCP), architectural decision records (ADR). Poor or misleading descriptions cause mis-selection and integration failures. Controlled mutation of description ``smells'' (e.g., functionality omission, accuracy misrepresentation) measures impact on selection probability; remediation can significantly improve correct tool choice~\cite{mcp-smell-arxiv}. Similarly, LLM-generated code must comply with ADRs; multi-model pipelines can detect violations with high accuracy, with RAG and human-in-the-loop closing gaps on implicit or deployment-oriented knowledge~\cite{adr-violation,mission-critical-rag}.

%-------------------------------------------------------------------------------
\section{Integration, Metrics, and Governance}
\label{sec:integration}
%-------------------------------------------------------------------------------

\subsection{Feedback Loop Between Reconstruction and Mutation}

Reconstruction consumes \emph{attack traces} and \emph{failure modes}; mutation produces them in a controlled way. A mutator (exemplar-guided) generates adversarial variants that stress the current policy; the reconstruction layer (SISF, context-construction defenses, hybrid judge ensemble) reacts and updates; the meta-cognitive monitor can trigger adversarial reset when cognitive instability is detected. The mutator then re-evaluates on new variants. Round scheduling and variant dispersion protect benchmark assets~\cite{ControlledMutation_SafetyEval}. Hygiene gates and \emph{mandatory} human calibration---a fixed proportion of high-confidence judge and policy decisions randomly audited---quantify and correct ensemble FPR~\cite{tc260-governance}.

\subsection{Metrics}

We align with a multi-dimensional protocol: (i) \textbf{Security}: RSR $\rightarrow$ 0 under mutation and scheduling; attacker AUROC/AUPRC $\approx$ random; decoy-hit rate and decoy-hit density (terminations only on high-fidelity extraction) monitored. (ii) \textbf{Fidelity}: Eval consistency on canonical vs.\ variant set kept small. (iii) \textbf{Reconstruction efficacy}: ASR reduction; judge-ensemble FPR via mandatory human calibration. (iv) \textbf{Mutation quality}: Naturalness, BERTScore, oracle consistency; exemplar adherence. (v) \textbf{Operational cost}: \textbf{End-to-end mutation latency} (time from failure signal to next admissible variant) and \textbf{token overhead ratio} (extra tokens for exemplar retrieval and injection vs.\ baseline) are mandatory metrics to ensure the framework is economically scalable at machine speed~\cite{ControlledMutation_SafetyEval}. (vi) \textbf{Audit}: Lineage and evidence chains for release gating (DUC/DUA-E style)~\cite{dual-use-attestations,ControlledMutation_SafetyEval}.

\subsection{Governance and Attestation}

Unified reconstruction and mutation support machine-executable compliance: evidence chains link each evaluation to mutator steps, scheduler state, and policy version. Release gates block deployment when critical variant families fail; diagnostics feed back to development. This aligns with dual-use attestations and pipeline verification standards (e.g., OWASP SPVS) and turns regulatory burden into a repeatable engineering process~\cite{ControlledMutation_SafetyEval}.

%-------------------------------------------------------------------------------
\section{Related Work}
%-------------------------------------------------------------------------------

LLM data extraction and reconstruction are well studied~\cite{llm-reconstruction-emergent,split-ft-attack,active-reconstruction}; we mitigate via structural mutation and scheduling rather than noise alone. Our context-construction defenses (ID alignment, decoy injection) extend intent-based firewalls~\cite{agentos-arxiv} to resist task reframing; our meta-cognitive monitor and exemplar-guided mutation harden against psychological manipulation; our hybrid judge ensemble and mandatory human calibration address single-judge bias. JADE and linguistics-based platforms~\cite{jade-linguistics,jade-db} inform mutation strategy; TC260 and attestation frameworks~\cite{tc260-governance,tc260-training-data,dual-use-attestations} underpin governance. SISF~\cite{sisf-arxiv}, AgentOS~\cite{agentos-arxiv}, ADM-ES~\cite{adm-es-arxiv}, CSE~\cite{cse-arxiv}, and epistemic grounding~\cite{epistemic-grounding} are building blocks we unify and harden; ControlledMutation\_SafetyEval~\cite{ControlledMutation_SafetyEval} provides the Mutator + Data Factory design.

%-------------------------------------------------------------------------------
\section{Conclusion}
%-------------------------------------------------------------------------------

The convergence of \emph{architectural reconstruction} and \emph{controlled mutation} defines a defensible path for advanced LLM safety. We have introduced \emph{deep architectural refactoring} to address hostile-examiner attack surfaces: (i) \textbf{context-construction defenses} (ID alignment, dynamic decoy injection, self-reminder) replace intent-only filtering and contaminate the attacker's operating space; (ii) \textbf{meta-cognitive monitoring} and \textbf{exemplar-guided mutation} anchor adaptation to golden exemplars and hierarchical experience memory, preventing value drift under psychological manipulation; (iii) a \textbf{hybrid judge ensemble} with distributional evaluation and mandatory human calibration removes single-judge bias from the feedback loop. The framework decouples safety from fragile LLM self-reflection: evolution is driven by decoys, exemplars, and a calibrated judge network, not by the model ``stepping on its own foot.'' As LLMs become infrastructure, this unified, hostile-examiner-hardened design is a necessary step toward scalable, auditable, and resilient safety defenses.

%-------------------------------------------------------------------------------
\section*{Availability}
%-------------------------------------------------------------------------------
Artifacts and extended proofs will be made available subject to institutional and export-control review.

%-------------------------------------------------------------------------------
\section*{Acknowledgments}
%-------------------------------------------------------------------------------
We thank our industry and standards colleagues for feedback on safety architecture and data governance.

%-------------------------------------------------------------------------------
\bibliographystyle{plain}
\bibliography{ArchitecturalReconstruction_ControlledMutation}

\end{document}
